\section{Conclusiones y Trabajo Futuro}
Tras desarrollar las distintas fases de ingeniería de características (añadiendo descomposiciones estacionales, medias móviles e indicadores binarios), queda claro que la calidad de estas variables es fundamental para mejorar métricas como el R$^2$ y el MAE. En este tipo de problemas, la preparación de los datos resulta más determinante que la arquitectura del propio modelo.

Como trabajo futuro, seguiremos iterando sobre los modelos \textbf{Random Forest} y \textbf{ARIMA}, ya que han demostrado ser los enfoques más robustos durante las pruebas. Cabe destacar que, aunque el modelo \textbf{ARIMA} presenta una alta sensibilidad inicial (con el dataset original con menos de 1000 filas) y requiere un histórico denso para ser preciso, su rendimiento mejora notablemente a medida que aumenta el volumen de datos. Esperamos que, con bases de datos más extensas, este modelo se consolide como la opción más fiable para el pronóstico a largo plazo.
